% ============================================================
% CAPITOLO 35 - LA TECNICA DELLA MODULAZIONE
% ============================================================
\chapter{La tecnica della modulazione}
\label{chap:tecnica_modulazione}

Il Capitolo \ref{chap:modulazione} ha dato la teoria della modulazione ai toni vicini. Qui si passa alla tecnica: quali settime fanno da pivot, come scegliere il pivot e confermare la nuova tonalità, come trattare le alterazioni e come modulare in fretta.

\section{Le specie di settima come strumenti modulanti}
\label{sec:specie_modulazione}

Le settime (Capitolo \ref{chap:sette_specie}) ampliano le possibilità modulanti rispetto alle sole triadi.

\subsection{La terza specie: il pivot per eccellenza}
\label{subsec:terza_specie_pivot}

La settima semidiminuita (terza specie, Capitolo \ref{chap:vii7_ii7}) ha una natura doppia: è \sigla{VIIø}{7}{} nel modo maggiore e \sigla{IIø}{7}{} nel minore relativo.

\begin{regola}[La terza specie come accordo pivot]
	\label{reg:terza_specie_pivot}
	La settima semidiminuita è il pivot ideale tra tonalità relative:
	\begin{itemize}
		\item \sigla{VIIø}{7}{} di Do maggiore $=$ \sigla{IIø}{7}{} di La minore;
		\item è lo stesso accordo (Si-Re-Fa-La) con due funzioni: dominante in Do M, sottodominante in La m.
	\end{itemize}
	L'ambiguità funzionale permette modulazioni molto fluide tra tonalità relative.
\end{regola}

\begin{esempio}[Modulazione con \sigla{VIIø}{7}{} $=$ \sigla{IIø}{7}{} come pivot]
	\label{ex:mod_terza_specie}
	L'accordo Si-Re-Fa-La funge da cerniera: è sentito come \sigla{VIIø}{7}{} (dominante) in Do maggiore, poi reinterpretato come \sigla{IIø}{7}{} (sottodominante) in La minore.
	\musica{es_cap16_04}
\end{esempio}

\subsection{La seconda specie: versatilità modulante}
\label{subsec:seconda_specie_pivot}

La settima minore (seconda specie, Capitolo \ref{chap:seconda_specie}) compare su molti gradi di scale diverse ed è molto versatile.

\begin{regola}[La seconda specie come pivot tra tonalità maggiori]
	\label{reg:seconda_specie_pivot}
	Esempi di accordi di seconda specie usabili come pivot:
	\begin{itemize}
		\item Re-Fa-La-Do: \sigla{II}{7}{} in Do M $=$ \sigla{VI}{7}{} in Fa M $=$ \sigla{IV}{7}{} in La m;
		\item Mi-Sol-Si-Re: \sigla{III}{7}{} in Do M $=$ \sigla{VI}{7}{} in Sol M $=$ \sigla{II}{7}{} in Re M.
	\end{itemize}
	La seconda specie permette di modulare ai toni vicini e di concatenare più modulazioni.
\end{regola}

\begin{esempio}[Modulazione Do M $\rightarrow$ Fa M con \sigla{II}{7}{} $=$ \sigla{VI}{7}{}]
	\label{ex:mod_seconda_specie}
	\musica{es_cap16_05}
\end{esempio}

\subsection{La quinta specie: modulazioni multiple}
\label{subsec:quinta_specie_modulazione}

La settima diminuita (Sezione \ref{sec:quinta_specie_dettaglio}) è simmetrica e può risolvere su quattro toniche diverse.

\begin{regola}[La quinta specie per modulazioni multiple]
	\label{reg:quinta_specie_modulazioni}
	Una stessa settima diminuita può essere reinterpretata come sensibile di quattro tonalità, distanti una terza minore l'una dall'altra. Si aprono così modulazioni ai toni vicini (quando la nuova tonica dista una terza minore) e ai toni lontani, che questo manuale non tratta.
\end{regola}

La quinta specie chiede molta attenzione nella preparazione e nella conferma della nuova tonalità, per non creare ambiguità.

\section{Scegliere il pivot}
\label{subsec:scelta_pivot}

\begin{regola}[Criteri per la scelta del pivot]
	\label{reg:scelta_pivot}
	Un buon pivot soddisfa questi requisiti:
	\begin{enumerate}
		\item \textbf{esiste in entrambe le tonalità} con le stesse note (o con alterazioni compatibili);
		\item \textbf{ha una funzione forte nella nuova tonalità:} meglio II, IV o V grado che III o VI;
		\item \textbf{permette una condotta fluida} dalla tonalità vecchia alla nuova;
		\item \textbf{non crea ambiguità premature:} è sentito come parte della tonalità vecchia finché non viene reinterpretato.
	\end{enumerate}
\end{regola}

\section{Confermare la nuova tonalità}
\label{subsec:conferma_tonalita}

\begin{definizione}[Conferma della tonalità]
	\label{def:conferma_tonalita}
	La \textbf{conferma della tonalità} è la presenza degli accordi che, insieme, identificano senza ambiguità la tonalità e il modo. Il metodo più diretto è la successione \textbf{IV-V-I} (o II-V-I): i tre accordi contengono \textit{tutte} le note della scala.
\end{definizione}

\begin{regola}[Importanza della cadenza di conferma]
	\label{reg:conferma_cadenza}
	Una modulazione non è completa senza una cadenza convincente nella nuova tonalità. Le opzioni più efficaci sono:
	\begin{itemize}
		\item \textbf{cadenza perfetta V--I:} la più forte e inequivocabile;
		\item \textbf{cadenza perfetta preceduta dalla sottodominante (II--V--I o IV--V--I):} rafforza ulteriormente la nuova tonalità;
		\item \textbf{ripetizione della cadenza:} nel dubbio, ripetere la progressione cadenzale elimina ogni ambiguità.
	\end{itemize}
	Non si interrompe la nuova tonalità subito dopo il pivot: la nuova tonica ha bisogno di tempo per stabilizzarsi.
\end{regola}

\section{Gestire le alterazioni}
\label{subsec:gestione_alterazioni}

\begin{regola}[Introduzione delle alterazioni]
	\label{reg:introduzione_alterazioni}
	Quando si modula a una tonalità con un diesis o un bemolle in più:
	\begin{itemize}
		\item \textbf{introdurre l'alterazione gradualmente:} compare di solito nella sensibile dell'accordo di dominante della nuova tonalità;
		\item \textbf{evitare salti cromatici bruschi:} se possibile l'alterazione appare in una voce interna e non nel soprano;
		\item \textbf{preparare l'orecchio:} il pivot dovrebbe già contenere alcune delle note «nuove» della scala di arrivo.
	\end{itemize}
\end{regola}

\begin{esempio}[Gestione del Fa$\sharp$ nella modulazione Do M $\rightarrow$ Sol M]
	\label{ex:gestione_fad}
	Il Fa$\sharp$ (sensibile di Sol maggiore) appare solo nell'accordo di dominante (Re-Fa$\sharp$-La-Do).
	\musica{es_cap16_06}
\end{esempio}

\section{Schemi di modulazione}
\label{sec:esempi_completi}

Le tre fasi si possono riassumere in uno schema di accordi, senza realizzazione a quattro voci.

\begin{center}
\small
\begin{tabularx}{\linewidth}{lLLL}
\toprule
\textbf{Percorso} & \textbf{Affermazione} & \textbf{Pivot} & \textbf{Conferma} \\
\midrule
Do M $\rightarrow$ La m & I (Do M) & II di Do M $=$ IV di La m & V--i in La m \\
Do M $\rightarrow$ Sol M & I (Do M) & I di Do M $=$ IV di Sol M & V--I in Sol M \\
Do M $\rightarrow$ Fa M & I (Do M) & I di Do M $=$ V di Fa M & V--I in Fa M \\
Do M $\rightarrow$ Mi m & I (Do M) & VI di Do M (La m) $=$ IV di Mi m & \sigla{IIø}{7}{}--V--i in Mi m \\
La m $\rightarrow$ Do M & i--III (La m) & VI di La m (Fa M) $=$ IV di Do M & IV--V--I in Do M \\
\bottomrule
\end{tabularx}
\end{center}

\section{Modulazione rapida}
\label{sec:regole_pratiche_modulazione}

Esistono procedimenti cromatici ricorrenti con cui da una tonalità maggiore si modula in fretta verso i toni vicini.

\begin{regola}[Modulazione al IV grado con la sensibile abbassata]
	\label{reg:mod_iv_sensibile_abbassata}
	In una tonalità maggiore, abbassando di un semitono la \textbf{sensibile} (VII grado) si ottiene la \textbf{settima} dell'accordo costruito sulla tonica. Questo accordo (\sigla{I}{7}{}) assume funzione di dominante della sottodominante e permette una modulazione immediata al \textbf{IV grado}. In Do maggiore il Si diventa \Sibemolle{}, l'accordo di Do diventa Do con settima e risolve su Fa maggiore.
\end{regola}

\begin{regola}[Modulazione al V grado con il IV grado innalzato]
	\label{reg:mod_v_quarto_innalzato}
	Per modulare al \textbf{V grado} (dominante) si innalza di un semitono il \textbf{IV grado} della tonalità di partenza. La nota diventa la \textbf{nuova sensibile} del V grado e trasforma l'accordo sul II grado in un accordo di dominante della dominante (\textbf{V/V}); sul IV grado alzato si ottiene invece la settima di sensibile della dominante (Fa$\sharp$-La-Do-Mi). In Do maggiore il Fa diventa Fa$\sharp$, sensibile di Sol maggiore.
\end{regola}

\begin{regola}[Modulazione al VI grado con il V grado innalzato]
	\label{reg:mod_vi_quinto_innalzato}
	Per modulare al \textbf{VI grado} (relativa minore della tonica) si innalza di un semitono il \textbf{V grado}. La nota diventa la \textbf{nuova sensibile} del VI grado e crea una dominante (\textbf{V/VI}) che risolve sulla nuova tonica minore. In Do maggiore il Sol diventa Sol$\sharp$, sensibile di La minore.
\end{regola}

\begin{regola}[Modulazione ai relativi minori degli altri gradi]
	\label{reg:mod_altri_minori}
	Si può modulare ai relativi minori degli altri gradi principali alterando i gradi della scala per ottenere la nuova sensibile:
	\begin{itemize}
		\item \textbf{al III grado (relativa del V):} si innalzano il \textbf{II} e il \textbf{IV grado} (in Do M: Fa$\sharp$ e Re$\sharp$ per modulare a Mi minore);
		\item \textbf{al II grado (relativa del IV):} si innalza il \textbf{I grado} (in Do M: Do$\sharp$ per modulare a Re minore).
	\end{itemize}
\end{regola}
